%! Author = sriadityavedantam
%! Date = 3/30/26

\section{Irreducibility}

\begin{proposition}[Polynomial Euclid's Lemma]{
    Suppose $p(x)\in\f[x]$ is irreducible and $b(x)\in\f[x]$ such that $p(x)\nmid b(x)$.
    Then
    \[
        \gcd(p(x),b(x))=1.
    \]
}
    Let $d(x)=\gcd(p(x),b(x))$.
    Then
    \[
        d(x)\mid p(x), \qquad d(x)\mid b(x).
    \]
    Since $p(x)$ is irreducible, every divisor of $p(x)$ is either a unit or an associate of $p(x)$.
    Because gcds are taken to be monic, $d(x)$ is monic.
    Hence either $d(x)=1$, or $d(x)=cp(x)$ for some nonzero $c\in\f$.
    If $d(x)=cp(x)$, then since $d(x)\mid b(x)$, it follows that $p(x)\mid b(x)$, which is a contradiction.
    Therefore
    \[
        \gcd(p(x),b(x))=1.
    \]
\end{proposition}

\begin{corollary}{
    If $p(x)\mid a_1(x)\cdots a_n(x)$, where $a_i(x)\in\f[x]$ and $p(x)$ is irreducible, then $p(x)\mid a_i(x)$ for some $i$.
}
    Show this by induction on $n$.
\end{corollary}

\begin{theorem}{
    Suppose $a(x)\in\f[x]$.
    Then $a(x)$ has a factorization into irreducible polynomials.
    This factorization is unique up to order and associates.
}
Use strong induction on $\deg(a(x))$ for existence.

For uniqueness, suppose
\begin{align*}
    a(x) &= p_1(x)\hdots p_r(x)\\
    &= q_1(x)\hdots q_s(x),
\end{align*}
where each $p_i(x)$ and $q_j(x)$ is irreducible.
Then
\[
    p_1(x)\mid q_1(x)\cdots q_s(x).
\]
By the previous corollary,
\[
    p_1(x)\mid q_j(x)
\]
for some $j$.
Without loss of generality, after rearranging, assume $j=1$.
Since $q_1(x)$ is irreducible, $p_1(x)$ and $q_1(x)$ are associates.
Proceed by induction on $\min\{r,s\}$.
\end{theorem}

\begin{lemma}[Irreducible Degrees]{
    Degree $1$ polynomials are irreducible.
    If a degree $2$ polynomial is reducible, then it is a product of linear polynomials.
}
\end{lemma}

\begin{lemma}[Freshman's Dream]{
    In $\z_2$,
    \[
        (x+1)^2=x^2+1.
    \]
}
\end{lemma}

\begin{proposition}{
    If $f(x)$ is irreducible in $\f[x]$, then so is every associate of $f(x)$.
}
\end{proposition}

Take note that for the equation $x^2+ax+b$, there are $p$ choices for each of $a,b$ in $\zp$, which means $p^2$ total monic choices for this polynomial.
More generally, the number of monic polynomials of degree $n$ in $\zp[x]$ is $p^n$.
The total number of polynomials of degree $n$ is
\[
    (p-1)p^n.
\]
\begin{example}{
    Prove that $x^2-2$ is irreducible in $\Q[x]$.
}
\end{example}

\begin{proof}
    We have
    \[
        x^2-2=\left(x-\sqrt{2}\right)\left(x+\sqrt{2}\right)\in\r[x].
    \]
    If $x^2-2$ factored in $\Q[x]$, then since it has degree $2$, it would factor into linear polynomials over $\Q$.
    That would force it to have a rational root.
    But its roots are $\pm\sqrt{2}$, and $\sqrt{2}\notin\q$.
    Therefore $x^2-2$ is irreducible in $\Q[x]$.
\end{proof}

Let $\f$ be a field.
Take $f(x)\in\f[x]$.
There is a corresponding polynomial function $\f\to\f$ also denoted by $f(x)$.

\begin{theorem}[Factor Theorem]{
    Let $f(x)\in\f[x]$ and $a\in\f$.
    If $f(a)=0$, then $(x-a)$ is a factor of $f(x)$.
    That is,
    \[
        f(x)=g(x)(x-a)
    \]
    for some $g(x)\in\f[x]$.
}
\end{theorem}

\begin{example}{
    (a) Show that $x^2+2$ is irreducible in $\mathbb{Z}_5[x]$.
}
\end{example}

\begin{proof}
    We will do a proof by contradiction.
    Suppose $x^2+2$ is reducible in $\mathbb{Z}_5[x]$.
    Since it has degree $2$, it must factor into linear polynomials:
    \[
        x^2+2=(x+a)(x+b)
    \]
    for some $a,b\in\mathbb{Z}_5$.
    Expanding gives
    \[
        (x+a)(x+b)=x^2+(a+b)x+ab.
    \]
    Since there is no degree $1$ term in $x^2+2$, we must have
    \[
        a+b=0,
    \]
    so $b=-a$.
    Then
    \[
        (x+a)(x-a)=x^2-a^2.
    \]
    Thus we would need
    \[
        -a^2=2
    \]
    in $\mathbb{Z}_5$, equivalently
    \[
        a^2=3.
    \]
    But the squares in $\mathbb{Z}_5$ are
    \[
        0^2=0,\quad 1^2=1,\quad 2^2=4,\quad 3^2=4,\quad 4^2=1.
    \]
    So $3$ is not a square in $\mathbb{Z}_5$.
    This is a contradiction.
    Therefore $x^2+2$ is irreducible in $\mathbb{Z}_5[x]$.
\end{proof}

(b) Factor $x^4-4$ as a product of irreducible in $\mathbb{Z}_5[x]$.

We have
\[
    x^4-4=(x^2+2)(x^2-2).
\]
Now in $\mathbb{Z}_5$, we have $-2\equiv 3$.
Since $4$ is a square in $\mathbb{Z}_5$, $x^2-2=x^2+3$ is reducible.
Indeed,
\[
    x^2-2=x^2+3=(x-1)(x+1).
\]
Also, from part (a), $x^2+2$ is irreducible.
Thus
\[
    x^4-4=(x^2+2)(x-1)(x+1)
\]
is a factorization into irreducible in $\mathbb{Z}_5[x]$.

\begin{theorem}[Remainder Theorem]{
    Let $f(x)\in\f[x]$ and $a\in\f$.
    Then
    \[
        f(x)=g(x)(x-a)+r(x)
    \]
    for some $g(x)\in\f[x]$, where $r(x)$ is a constant polynomial.
}
\end{theorem}

\begin{proof}
    By the division algorithm,
    \[
        f(x)=g(x)(x-a)+r(x),
    \]
    where $r(x)=0$ or
    \[
        \deg(r(x))<\deg(x-a)=1.
    \]
    Thus $r(x)$ must be a constant polynomial or $0$.
\end{proof}

\begin{proof}
    By the remainder theorem,
    \[
        f(x)=g(x)(x-a)+r
    \]
    where $r$ is a constant.
    Evaluating at $x=a$ gives
    \begin{align*}
        f(a) &= g(a)(a-a)+r\\
        &= r.
    \end{align*}
    So if $f(a)=0$, then $r=0$.
    Hence
    \[
        f(x)=g(x)(x-a).
    \]
\end{proof}

\begin{definition}[Roots]
    $a$ is a root of $f(x)$ if
    \[
        f(a)=0.
    \]
\end{definition}

\begin{corollary}[of Factor Theorem]{
    Suppose $f(x)\in\f[x]$ has $\deg f(x)=n$.
    Then $f(x)$ has at most $n$ different roots.
}
    By induction on $\deg f(x)$.
    If $\deg f(x)=0$, then $f$ is a nonzero constant and has no roots.
    If $\deg f(x)=1$, then
    \[
        f(x)=a_{1}x+a_0,
        \qquad a_1\neq 0.
    \]
    It has exactly one root, namely
    \[
        x=\frac{-a_0}{a_1}.
    \]
    Assume true for polynomials of degree $n-1$.
    Now let $\deg f(x)=n$ and suppose $a$ is a root of $f(x)$.
    Then by the factor theorem,
    \[
        f(x)=(x-a)g(x)
    \]
    for some $g(x)$ with $\deg g(x)=n-1$.
    If $b\neq a$ is another root of $f(x)$, then
    \[
        0=f(b)=(b-a)g(b).
    \]
    Since $b-a\neq 0$, we get $g(b)=0$.
    Thus every root of $f(x)$ other than $a$ is a root of $g(x)$.
    By the induction hypothesis, there are at most $n-1$ such roots.
    So the number of roots of $f(x)$ is at most
    \[
        1+(n-1)=n.
    \]
\end{corollary}

If $f(x)\in\Q[x]$, then the rational root test tells us if $f(x)$ has a linear factor.

\begin{definition}[Rational Root Test]
    Suppose
    \[
        f(x)=a_{n}x^n+\hdots+a_0\in\z[x].
    \]
    If $\frac{r}{s}\in\q$ in lowest terms is a root of $f(x)$, then
    \[
        r\mid a_0
        \qquad\text{and}\qquad
        s\mid a_n.
    \]
    Since $a_0$ and $a_n$ have finitely many divisors, there are only finitely many possibilities to check.
\end{definition}

For example, $2x^3-x^2+1$ is irreducible by the rational root test, since the only possible rational roots are
\[
    \pm 1,\ \pm \frac{1}{2}.
\]
After checking all possibilities, none of them are roots.
Suppose $f(x)\in\zx$ and $g(x),h(x)\in\qx$.
Then there exist $\alpha,\beta\in\q$ such that
\[
    f(x)=(\alpha g(x))(\beta h(x))\in\zx.
\]
Also, if $f(x)\in\qx$, then there is a $c\in\q$ such that $cf(x)\in\zx$.
The rational root test tells us about linear factors, which suffices to show irreducibility for degrees $2$ and $3$, but not in higher degrees.
This builds the foundation for the following theorem.

\begin{theorem}[Gauss's Lemma of Irreducibility]{
    Suppose $f(x)\in\zx$.
    If $f(x)$ is irreducible in $\zx$, then $f(x)$ is irreducible in $\qx$.
}
\end{theorem}

One may ask whether the converse is possible given these assumptions.
I claim not always.
It is possible when primitivity is included.
This course will not go deeply into primitivity, but here is the definition.

\begin{definition}[Primitivity]
    A polynomial $p(x)$ with integer coefficients is called primitive if and only if the gcd of all its coefficients is $1$.
\end{definition}

If this is also true, then and only then do we get the full biconditional statement.
This was a whole block of assumptions to unfold before displaying the if-then statement of Gauss's lemma of irreducibility.
But let us look at an example of how to apply this.

Let
\[
    f(x)=6x^2-5x+1.
\]
Then
\[
    f(x)=\left(x-\frac{1}{2}\right)(6x-2),
\]
so
\[
    f\left(\frac{1}{2}\right)=0.
\]
Thus we have shown a root in $\qx$, which demonstrates that it is reducible.
But we can also write this in the form of integer factors:
\[
    f(x)=(2x-1)(3x-1)\in\zx.
\]

\begin{lemma}[Introductory Lemma]{
    Suppose $f(x),g(x),h(x)\in\zx$ where
    \[
        f(x)=g(x)h(x).
    \]
    Let $p$ be prime such that $p$ divides every coefficient of $f(x)$.
    Then either $p$ divides every coefficient of $g(x)$ or $p$ divides every coefficient of $h(x)$.
}
\end{lemma}

\begin{proof}
    Suppose
    \[
        f(x)=g(x)h(x),
    \]
    with
    \[
        f(x)=a_0+a_{1}x+\hdots,\quad g(x)=b_0+b_{1}x+\hdots,\quad h(x)=c_0+c_{1}x+\hdots.
    \]
    Then
    \[
        a_0=b_{0}c_0.
    \]
    Since $p\mid a_0$, by Euclid's lemma we get
    \[
        p\mid b_0
        \qquad\text{or}\qquad
        p\mid c_0.
    \]
    Assume without loss of generality that $p\nmid c_0$.
    We will show by induction that $p\mid b_i$ for all $i$.
    Now
    \[
        a_1=b_{0}c_1+b_{1}c_0.
    \]
    Since $p\mid a_1$ and $p\mid b_{0}c_1$, it follows that $p\mid b_{1}c_0$.
    Because $p\nmid c_0$, we conclude $p\mid b_1$.
    Continuing this argument inductively shows that $p$ divides every coefficient of $g(x)$.
\end{proof}

\begin{theorem}[Eisenstein's Theorem of Irreducibility]{
    Suppose
    \[
        f(x)=a_{n}x^n+\hdots+a_0\in\zx.
    \]
    Suppose $p\nmid a_n$, $p\mid a_i$ for all $i<n$, and $p^2\nmid a_0$.
    Then $f(x)$ is irreducible in $\qx$.
}
Suppose $f(x)$ is reducible in $\qx$.
Then by Gauss's lemma, $f(x)$ is reducible in $\zx$.
So
\[
    f(x)=g(x)h(x),
\]
where $g(x),h(x)\in\zx$ and
\[
    \deg g(x),\deg h(x)<\deg f(x)=n.
\]
Using the introductory lemma and tracking the divisibility of coefficients by $p$, one concludes that $p$ must divide every coefficient of one factor.
Then $p^2$ divides the constant term $a_0$, contradicting the hypothesis.
Therefore $f(x)$ is irreducible in $\qx$.
\end{theorem}

\begin{lemma}{
    Linear polynomials are not reducible.
}
    A linear polynomial cannot be written as a product of two non-unit polynomials, since degrees add under multiplication and both non-unit factors would have positive degree.
\end{lemma}

Let
\[
    f(x)=2x^4+15x^3+30x^2+60x-21.
\]
We have
\[
    3\nmid 2,\qquad 3\mid 15,30,60,21,\qquad 9\nmid 21.
\]
So $f(x)$ is irreducible by Eisenstein.

\begin{theorem}[Reduction mod $p$]{
    Let $f(x)\in\zx$.
    Let $p\nmid a_n$, where $a_n$ is the leading coefficient of $f(x)$.
    Consider
    \[
        \overline{f(x)}=\overline{a_n}x^n+\hdots+\overline{a_0},
    \]
    where $\overline{a_i}$ is the congruence class of $a_i\mod p$.
    If $\overline{f(x)}$ is irreducible in $\zpx$, then $f(x)$ is irreducible in $\zx$ and $\qx$.
    The converse is not true.
}
Suppose $f(x)$ is reducible in $\zx$.
Then
\[
    f(x)=g(x)h(x)
\]
for some $g(x),h(x)\in\zx$ of smaller positive degree.
Reducing mod $p$ gives
\[
    \overline{f(x)}=\overline{g(x)}\,\overline{h(x)}.
\]
Since $p\nmid a_n$, the degree of $\overline{f(x)}$ is still $n$, so both $\overline{g(x)}$ and $\overline{h(x)}$ have smaller degree than $\overline{f(x)}$.
This contradicts the irreducibility of $\overline{f(x)}$ in $\zpx$.
So $f(x)$ is irreducible in $\zx$.
Then by Gauss's lemma, $f(x)$ is irreducible in $\qx$.
\end{theorem}

Let
\[
    f(x)=x^4+3x^3+6x^2+1\in\Q[x].
\]
Try $p=2$.
Then
\[
    \overline{f(x)}=x^4+x^3+1.
\]
It has no roots in $\z_2$, so it has no linear factors.
Hence it is irreducible in $\z_2[x]$, and therefore $f(x)$ is irreducible in $\Q[x]$.

\begin{theorem}[Fundamental Theorem of Algebra]{
    If $f(x)\in\cx$, then the following are equivalent:
    \begin{enumerate}
        \item $f(x)$ is irreducible;
        \item $f(x)$ is linear;
        \item every non-constant polynomial in $\cx$ can be factored as a product of linear factors;
        \item every non-constant polynomial in $\cx$ has a root.
    \end{enumerate}
}
\end{theorem}

For example,
\[
    f(x)=x^2+1\in\cx
\]
has complex roots $\pm i$, and
\[
    f(x)=(x+i)(x-i).
\]
Let $\theta=\frac{2\pi}{3},\frac{4\pi}{3}$.
Then
\begin{align*}
    e^{\theta i} &= \cos\theta+i\sin\theta,\\
    \left(e^{\theta i}\right)^3 &= \cos 3\theta+i\sin 3\theta\\
    &= \cos 2\pi+i\sin 2\pi\\
    &= \cos 4\pi+i\sin 4\pi\\
    &= 1.
\end{align*}
Roots of $x^n-1$ are
\[
    e^{\theta i},e^{2\theta i},e^{3\theta i},\hdots,e^{(n-1)\theta i}.
\]

\begin{proposition}{
    Suppose $f(x)\in\rx$.
    Every irreducible polynomial in $\rx$ has degree $1$ or $2$.
}
    Consider $f(x)\in\rx$ as a polynomial in $\cx$.
    By the Fundamental Theorem of Algebra, $f(x)$ has a root in $\c$.
    If this root is real, then $f(x)$ has a linear factor.
    So assume the root is
    \[
        \omega=a+bi
    \]
    with $b\neq 0$.
    Then $\overline{\omega}=a-bi$ is also a root.
    Indeed, since the coefficients of $f(x)$ are real, complex conjugation fixes each coefficient, and thus
    \[
        f(\omega)=0 \implies f(\overline{\omega})=0.
    \]
    Therefore
    \[
        (x-\omega)(x-\overline{\omega})
    \]
    is a factor of $f(x)$.
    But
    \[
        (x-\omega)(x-\overline{\omega})=x^2-(\omega+\overline{\omega})x+\omega\overline{\omega},
    \]
    and
    \[
        \omega+\overline{\omega}\in\r,\qquad \omega\overline{\omega}\in\r.
    \]
    So this quadratic factor lies in $\rx$.
    Hence every irreducible polynomial in $\rx$ has degree $1$ or $2$.
\end{proposition}

\begin{example}
    Suppose $f(x)\in\rx$ and has degree $3$.
    By the Intermediate Value Theorem, there exists $c\in\r$ such that $f(c)=0$.
    So by the factor theorem, $(x-c)$ is a factor of $f(x)\in\rx$.
\end{example}

\begin{lemma}{
    Now let $\phi(a+bi)=a-bi$.
    Then $\phi:\c\to\c$ is an isomorphism.
}
\end{lemma}